% !TEX root = ../paper.tex
% Section: Introduction

\section{Introduction}
\label{sec:intro}

Two-layer linear networks are nonconvex in their weight factors, even when the
loss is a simple function of the predictor matrix.  Training moves in the
lifted space of incoming and outgoing neuron weights, whose second moments act
as a time-dependent preconditioner on the predictor.  These moments may in
principle become degenerate and slow down the dynamics, and it is not clear that the dynamics necessarily converge to a globally optimal predictor.
We consider this question in the rich mean-field scaling, in which both sides of each rank-one neuron evolve, and
smooth predictor losses satisfying a Polyak--\L{}ojasiewicz (PL) inequality, which subsume the standard quadratic loss.
% Unlike lazy or neural tangent kernel scalings, the hidden representation is not
% frozen and its conditioning must be controlled globally in time.

Central to our analysis are width-independent conservation laws for
the weight covariance: the signed covariance spectrum, relative to the input-output
block split, supplies conserved supporting inequalities and lower bounds on each hidden-block. The sum of these lower bounds conditions the predictor dynamics: a positive sum
is sufficient to guarantee global convergence of the predictor loss with an explicit rate. This analysis covers every positive-definite initialization and some singular ones.
For initial variance $\sigma^2 \Id$, the loss decays with rate at least $4\sigma^2\kappa$, where $\kappa$ is the PL constant.

The covariance viewpoint also makes the result robust to finite-width
sampling and factor gradient descent. For deep linear ResNets, the same
conservation argument applied layerwise yields linear convergence
conditional on a uniform bound on the residual path. With heavy-ball
momentum, closure is lifted to the joint phase covariance of positions
and velocities. Conformally transported spectral laws describe its
geometry, while energy dissipation gives an explicit initialization and
damping certificate for linear convergence. Special values of the
conserved quantities further permit lower-dimensional intrinsic dynamics;
these reductions are described in Appendix~\ref{app:intrinsic}.

\paragraph{Related work}
% \paragraph{Mean-field flows and feature learning.}
Mean-field analyses describe the collective evolution of neuron weight
distributions in wide two-layer networks~\citep{ChizatBach2018,MeiMontanariNguyen2018,RotskoffVandenEijnden2018}.
These formulations treat broad nonlinear feature maps and often prove asymptotic global optimality or propagation-of-chaos type statements under support, homogeneity, noise, or regularity assumptions.
In the present linear setting, exact second-moment closure connects
training to Bures--Wasserstein covariance geometry~\citep{Takatsu2011,BhatiaJainLim2019},
making a sharp quantitative linear convergence rate available.
This closure distinguishes our analysis from the NTK~\citep{JacotGabrielHongler2018} and lazy training~\citep{ChizatOyallonBach2019} regime: both incoming and outgoing weights move, and the evolving representation is essential to the rate.

% \paragraph{Linear networks and conservation laws.}
Deep linear networks have served as a tractable model of the nonconvexity created by matrix products, with results on their landscapes, exact dynamics,
and convergence~\citep{SaxeMcClellandGanguli2014,Kawaguchi2016,LaurentVonBrecht2018,AroraCohenGolowichHu2019}.
Finite-width analyses of two-layer matrix factorization furnish convergence guarantees under small, asymmetric, or general initializations~\citep{YeDu2021,XuMinTarmounMalladaVidal2023}, while recent work treats discrete-time low-rank adaptation directly~\citep{MuKlabjan2025}.
Our result instead starts from a population covariance satisfying a positive \emph{support gap} condition, in particular any positive-definite covariance, and obtains an explicit global rate that survives independent finite-width sampling.
Related work identifies Gramian differences and balancedness as invariants of factor training~\citep{DuHuLee2018,TarmounFrancaHaeffeleVidal2021} and develops
a general differential framework for identifying conservation laws~\citep{MarcotteGribonvalPeyre2023}.
Our mean-field laws are functionals of the neuron distribution that induce
compatible particle laws at every width. The conserved spectrum of the
signed covariance yields hidden-block lower bounds without requiring
a special balancedness identity.

% \paragraph{Deep residual networks}
Residual networks~\citep{HeZhangRenSun2016Deep,HeZhangRenSun2016Identity} were introduced as architectures built from perturbations of the identity, and their infinite depth limit leads naturally to neural ODEs~\citep{HeZhangRenSun2016Deep,ChenRubanovaBettencourtDuvenaud2018}.
Mean-field formulations have established global optimality and convergence results under expressivity, regularity, or local PL conditions~\citep{LuMaLuLuYing2020,DingChenLiWright2021,Isobe2023,BarboniPeyreVialard2024COTResNets}.
Our linear analysis provides a complementary condition: uniform conditioning of the forward and backward residual products is sufficient for convergence.

% \paragraph{Heavy-ball momentum.}
Momentum turns first-order descent into an inertial, second-order iteration; the heavy-ball method of~\cite{Polyak1964} and Nesterov acceleration~\citep{Nesterov1983} have mature convergence theories, including continuous-time interpretations~\citep{SuBoydCandes2016} and robust-control analyses~\citep{LessardRechtPackard2016}.
For deep linear networks, \citet{WangLinAbernethy2021,LiuTaoPan2022}
establish acceleration results under initialization and overparameterization
conditions. Here the appropriate state variable is the phase covariance, over which the dynamics close.
We give an explicit energy-damping certificate for global convergence
with general position-velocity correlations.

\paragraph{Our contributions}
Section~\ref{sec:laws} identifies width-independent spectral conservation
laws and derives uniform hidden-block bounds from their support margins
(Propositions~\ref{prop:covariance-closure-laws} and~\ref{lem:margins}).
Their Wasserstein interpretation and local completeness among smooth covariance
laws common to all split losses are developed in
Appendices~\ref{app:mean-field} and~\ref{app:law-counting}.
Section~\ref{sec:convergence} proves global exponential convergence with an explicit rate under a smooth PL predictor loss from every initial covariance with positive
support gap (Theorem~\ref{thm:gf}).
The rate survives independent finite-width sampling at width proportional to the population covariance rank, which can be
as small as $\min(d_u,d_v)$ (Corollary~\ref{cor:finite-width} and
Remark~\ref{rem:rank-adaptive-width}).
We establish global linear convergence of factor
gradient descent under an explicit stepsize bound by controlling the
accumulated conservation error (Theorem~\ref{thm:gd}).
Section~\ref{sec:resnets} extends the layerwise argument to finite- and
continuous-depth linear ResNets under residual-path control
(Theorems~\ref{thm:resnet-finite} and~\ref{thm:resnet-continuous}), and
discusses why that control is not automatic (Remark~\ref{rem:resnet-path-escape}
in Appendix~\ref{app:resnets}).
Section~\ref{sec:momentum} shows that heavy-ball momentum dynamics closes on the phase covariance and identifies conformal conservation laws (Proposition~\ref{prop:phase-laws}).
Global exponential convergence is proven under an initial-energy
and damping condition, including a two-scale white specialization
(Theorem~\ref{thm:momentum} and Corollary~\ref{cor:two-scale-white-window}).

%%% Local Variables:
%%% mode: LaTeX
%%% TeX-master: "../paper"
%%% End:
